\newpage
\section{Introducción}
En las últimas décadas, el acceso a Internet ha pasado de ser un servicio complementario a convertirse en un recurso indispensable para el desarrollo de la educación, la economía, la investigación científica, las comunicaciones y múltiples actividades de la vida cotidiana. Sin embargo, millones de personas aún enfrentan dificultades para acceder a una conexión de calidad debido a limitaciones geográficas, económicas o de infraestructura. Frente a este escenario, el internet satelital ha cobrado una importancia creciente al ofrecer una alternativa capaz de extender la conectividad hacia regiones donde las redes terrestres tradicionales resultan insuficientes o inviables.
El desarrollo de esta tecnología ha experimentado una transformación acelerada gracias a los avances en la industria espacial, la reducción de los costos de lanzamiento, la miniaturización de los satélites y el despliegue de grandes constelaciones orbitales impulsadas por empresas como SpaceX, Amazon y otros actores del sector. Paralelamente, el internet satelital ha dejado de ser únicamente una solución para zonas aisladas y ha pasado a formar parte de una estrategia global que involucra aspectos tecnológicos, económicos, ambientales e incluso geopolíticos, al influir en la competencia internacional por el liderazgo en las telecomunicaciones espaciales y en el desarrollo de nuevas capacidades de conectividad.
En este contexto, el presente trabajo tiene como propósito analizar el internet satelital desde una perspectiva integral. Para ello, se desarrollan sus fundamentos técnicos, su arquitectura, principio de funcionamiento y evolución tecnológica, así como las principales características de las órbitas satelitales y las innovaciones más recientes en materia de conectividad. Asimismo, se examinan las aplicaciones actuales de esta tecnología, sus ventajas y limitaciones, los impactos ambientales derivados del creciente despliegue de satélites, la competencia tecnológica entre potencias como Estados Unidos y China, el papel de empresas privadas en la denominada nueva economía espacial y la situación del internet satelital en el Perú.

\section{Fundamentos del Internet satelital}
\subsection{Concepto y principio de funcionamiento}
\subsection{ Arquitectura del sistema de Internet satelital}
\subsection{ Importancia de la latencia como factor de rendimiento}
\section{Evolución tecnológica}
\subsection{Antecedentes de las comunicaciones satelitales}
\subsection{Evolución de los sistemas satelitales: comparación entre GEO, MEO y LEO}
\subsection{Direct-to-Cell: comunicación directa con dispositivos móviles}
